% Answers to Chapter 13, translated from sv/back/facit/13.tex.

\facitavsnitt{c13}{Chapter 13}{Derivatives, part II}
\begin{facit}

\svar{1.1} \sv{a} $f'(x) = -1.5x^2 + 3x + 4.5$ \sv{b} $x = -1$ and $x = 3$
\sv{c} Maximum at $x = 3$ ($f''(3) = -6 < 0$), minimum at $x = -1$ ($f''(-1) = 6 > 0$).
\sv{d} The maximum value is $f(3) = 10.5$ and the minimum value $f(-1) = -5.5$.
\sv{e} The curve has a valley at $(-1, -5.5)$ and a peak at $(3, 10.5)$, and crosses the $y$-axis
at $-3$.

\svar{1.2} \sv{a} $u'(t) = -0.6t^2 + 2.4t - 1.4625$
\sv{b} $t_1 = 0.75\,\text{s}$ and $t_2 = 3.25\,\text{s}$
\sv{c} $u''(0.75) = 1.5 > 0$: minimum; $u''(3.25) = -1.5 < 0$: maximum.
\sv{d} The minimum value is $u(0.75) \approx 2.49\,\text{V}$ and the maximum value
$u(3.25) \approx 4.06\,\text{V}$.
\sv{e} The curve starts at $3\,\text{V}$, has a valley at $(0.75, 2.49)$ and a peak at
$(3.25, 4.06)$, and then falls.

\svar{2.1} \sv{a} $i'(t) = 200\cos(50t)$
\sv{b} $50t = \pi/2 + k\pi$, so $t = \pi/100 + k\pi/50$,
$k = 0, 1, 2, \ldots$; the first two are $t_1 \approx 0.031\,\text{s}$ and
$t_2 \approx 0.094\,\text{s}$.
\sv{c} $i''(t) = -10\,000\sin(50t)$: $i''(t_1) = -10\,000 < 0$ gives a maximum,
$i''(t_2) = 10\,000 > 0$ a minimum.
\sv{d} $4\,\text{A}$ and $-4\,\text{A}$, which matches the amplitude $4\,\text{A}$.

\svar{2.2} \sv{a} $u'(t) = -4.8e^{-0.4t}$ \sv{b} $u'(2) \approx -2.16\,\text{V/s}$
\sv{c} $u''(2) \approx 0.86\,\text{V/s}^2 > 0$: the rate of change is increasing, that is, the
voltage is falling more and more slowly.

\svar{2.3} \sv{a} $x = (10^{0.2} - 1)/2 \approx 0.29$
\sv{b} $G'(x) = 20/\bigl((2x + 1)\ln 10\bigr) \approx 5.48$ dB per unit

\svar{3.1} \sv{a} $3\cos(3x)$ \sv{b} $-5\sin(5x)$ \sv{c} $4\cos(x)$ \sv{d} $4\sin(2x)$
\sv{e} $\cos(x) - \sin(x)$

\svar{3.2} \sv{a} $5e^{5x}$ \sv{b} $-6e^{-2x}$ \sv{c} $e^x + 2x$ \sv{d} $2^x \ln 2$
\sv{e} $-2.5e^{-x/4}$

\svar{3.3} \sv{a} $\frac{1}{x}$ \sv{b} $\frac{1}{x}$ \sv{c} $\frac{3}{x}$
\sv{d} $\frac{1}{x \ln 10}$ \sv{e} $\frac{1}{x} + e^x$

\svar{3.4} \sv{a} $6e^{3x} - 4\cos(4x)$ \sv{b} $2x + \frac{1}{x}$ \sv{c} $-5\sin(x) + 3e^{-x}$
\sv{d} $8\cos(2x) - 8\sin(2x)$

\svar{3.5} \sv{a} $2$ \sv{b} $0$ \sv{c} $0.5$ \sv{d} $0$ \sv{e} $-3e^{-1} \approx -1.10$

\svar{3.6} \sv{a} $u'(t) = 1000\pi\cos(100\pi t)$
\sv{b} $i_C(t) \approx 0.148\cos(100\pi t)\,\text{A}$ \sv{c} About $148\,\text{mA}$
\sv{d} At $t = 0$, when the voltage passes through zero and changes fastest.

\svar{3.7} \sv{a} $i'(t) = 400\cos(200t)$ \sv{b} $u_L(t) = 10\cos(200t)\,\text{V}$
\sv{c} $10\,\text{V}$ \sv{d} The voltage leads the current by $90^{\circ}$.

\svar{3.8} \sv{a} $u'(t) = -40e^{-2t}\,\text{V/s}$ \sv{b} $-40\,\text{V/s}$
\sv{c} $-40e^{-1} \approx -14.7\,\text{V/s}$
\sv{d} The derivative is negative, so the voltage is decreasing, fastest at the start.
\sv{e} $-u(t)/0.5 = -20e^{-2t}/0.5 = -40e^{-2t}$, the same expression as $u'(t)$ in a).

\svar{3.9} \sv{a} $f'(x) = e^{-2x}(1 - 2x)$ \sv{b} $x = 0.5$
\sv{c} Maximum: $f''(x) = e^{-2x}(4x - 4)$ gives $f''(0.5) = -2e^{-1} \approx -0.736 < 0$.
\sv{d} $f(0.5) = 0.5e^{-1} \approx 0.184$

\svar{3.10} \sv{a} $p'(t) = 100e^{-t}(1 - t)$ \sv{b} $t = 1\,\text{s}$
\sv{c} Maximum, since $p''(1) = -100e^{-1} < 0$.
\sv{d} $p(1) = 100e^{-1} \approx 36.8\,\text{W}$

\svar{3.11} \sv{a} $G'(x) = \frac{20}{x \ln 10}$ \sv{b} $\approx 8.69$ dB per unit
\sv{c} $\approx 0.869$ dB per unit
\sv{d} The derivative is ten times smaller at $x = 10$: the gain in dB grows more and more slowly,
since the logarithmic scale compresses large values.

\svar{3.12} \sv{a} $1$ \sv{b} $-1$ \sv{c} $y = 1 - x$
\sv{d} At $x = 1$, that is, after one time constant ($\tau = 1$): had the discharge continued as
fast as at the start, it would have been complete after the time $\tau$.

\svar{3.13} \sv{a} With $k = -1/\tau$, $u'(t) = -U_0e^{-t/\tau}/\tau = -u(t)/\tau$.
\sv{b} The rate of change is proportional to the present voltage: the lower the voltage, the more
slowly it falls.
\sv{c} $u'(0) = -5\,\text{V/s}$, $u'(\tau) \approx -1.84\,\text{V/s}$
\sv{d} $e^{-1} \approx 37\,\%$

\svar{3.14} \sv{a} False: the derivative is $\cos x$. \sv{b} False: it is $3e^{3x}$.
\sv{c} False: $\ln(5x) = \ln 5 + \ln x$, so the derivative is $\frac{1}{x}$.
\sv{d} True: $\cos(kx)$ has the derivative $-k\sin(kx)$, here with $k = 2$.
\sv{e} True: $e^{kx}$ has the derivative $ke^{kx}$, and here $k = 1$.
\sv{f} True: the constant can be moved outside, $\bigl(Cf(x)\bigr)' = Cf'(x)$.

\svar{3.15} \sv{a} The inner factor $3$ must be included: $3\cos(3x)$.
\sv{b} The exponent does not change; only a factor is added: $-2e^{-2x}$.
\sv{c} The constant inside the logarithm disappears: $\frac{1}{x}$.
\sv{d} The derivative of cosine is minus sine: $-4\sin(x)$.
\sv{e} The power rule does not apply when $x$ is in the exponent: $2^x \ln 2$.

\svar[Test yourself]{T} \sv{1} $f'(x) = 2e^{2x} + \frac{1}{x}$
\sv{2} $f'(\pi/2) = \cos(\pi/2) = 0$

\end{facit}
